\documentclass[11pt]{article}
\usepackage[review]{acl}
\usepackage{times}
\usepackage{latexsym}
\usepackage[T1]{fontenc}
\usepackage[utf8]{inputenc}
\usepackage{microtype}
\usepackage{graphicx}
\usepackage{booktabs}
\usepackage{amsmath}

% PAPER 1 (línea ANN/memoria). Esqueleto liviano — JM conoce mejor este trabajo, llenarlo.
\title{Dynamic Turn Embeddings for Conversational Memory Retrieval in Task-Oriented Dialogue}

\author{Juan Manuel Fern\'andez \\ UNSL \\ \texttt{...}}

\begin{document}
\maketitle

\begin{abstract}
\textit{[abstract — TODO]}
% Historia: recuperación de memoria conversacional vía búsqueda ANN sobre representaciones de
% turno. Comparamos representaciones estáticas, acumulativas-dinámicas-normalizadas y por EMA.
% Hallazgo: para RETRIEVAL, las memorias a mano (EMA/acumulativo) igualan o superan a encoders
% contextuales -> "contextualizar != recuperar".
\end{abstract}

\section{Introduction}
% Memoria conversacional; recuperación de situaciones por vecinos más cercanos (ANN). TODO.

\section{Related Work}
% Sentence/turn encoders, ANN/FAISS, memoria en diálogo. TODO.

\section{Method}
% Representaciones sobre e_t: Static, Cumulative (promedio normalizado), EMA(alpha).
% Índice ANN (FAISS). TODO.

\section{Evaluation}
% Métrica no-circular: MSS@10 con juez LLM (situación). act-match P@k (circular). TODO.

\section{Results}
% Hallazgo central: EMA/acumulativo >= contextual en retrieval (MSS). Tabla. TODO.

\section{Discussion}
% "Contextualiza != recupera": e_t (D2F) ya es un espacio retrieval-tuned; la memoria a mano se
% queda en él y gana; el BERT-de-turnos driftea (contextualiza) y pierde el ancla. TODO.

\section{Conclusion}
% TODO.

\section*{Limitations}
% TODO.

% \bibliography{references}  % TODO(JM): descomentar al agregar referencias (acl.sty fija el estilo)

\end{document}
